% k01 King's existence theorem on the cube
% local macros for the named statements of this section and the next (to be moved to preamble.tex)
\providecommand{\moving}[1]{\textcolor{orange!80!black}{[#1]}}
\providecommand{\Hrel}{\ensuremath{(\mathrm{H0rel})_{C}}}
\providecommand{\Kpos}{\ensuremath{(\mathrm{K}\text{-}\mathrm{pos})_{N}}}
\providecommand{\BALint}{\ensuremath{(\mathrm{BAL}_{5})[\mathrm{int}]}}
\providecommand{\RENorb}{\ensuremath{(\mathrm{REN}_{\mathrm{orb}})}}
\providecommand{\Kh}{\ensuremath{(\mathrm{K}_{h})}}
\providecommand{\Uplus}{\ensuremath{(\mathrm{U}^{+})}}
\providecommand{\Bh}{\ensuremath{(\mathrm{B}^{h})}}
\providecommand{\Rzero}{\ensuremath{(\mathrm{R}_{0})}}
\providecommand{\UBm}{\ensuremath{(\mathrm{UB}^{-})}}
\providecommand{\NEGinf}{\ensuremath{(\mathrm{NEG}_{\infty})}}
\providecommand{\Zk}{\ensuremath{Z_{k}}}
\providecommand{\Ik}[1]{\ensuremath{I^{(#1)}_{k}}}
\providecommand{\Ihat}[1]{\ensuremath{\hat I^{(#1)}_{k}}}
\providecommand{\Cl}{\operatorname{Cl}}

\section{King's existence theorem on the cube}\label{sec:k01}

The existence theorem for King's cube limits (\cref{hyp:s03:EC}), and its form with scalar sources that the Wick leg consumes (\cref{hyp:s03:W1}), are proved in this section and the next from estimates on King's renormalisation-group integrals on the cube $\CK$. The argument is the one King gives for his Theorem~2.1 \cite[Sect.~3.1]{K2} and \cite[Sect.~3.2]{K1}: the generating functional is approximated, up to a small-field remainder, by an integral of the effective action after $j$ renormalisation steps, and two spacings whose effective actions live on the same unit lattice are compared. We write it as a Cauchy property of the \emph{normalised} functional, in a ratio form in which every field-independent constant of the effective action cancels; this removes from King's route the convergence of the vacuum energies, which his printed argument uses and which is not available on the cube.

The inputs are of two kinds. Estimates taken from Ba{\l}aban's block-factor-$5$ analysis \cite{Ba1,Ba2,Ba3,Ba4} on tori, and their transfer to King's cube, are stated as propositions, each followed by a description of its content; their proofs are collected in an appendix. Three of these statements are used at several places and are therefore stated once, with a precise formulation, in \cref{ssec:k01:open} (\BALint, \RENorb) and \cref{ssec:k02:setting} (\Kh). Everything else in the two sections is proved in place.

\subsection{Renormalisation-group integrals on King's cube}\label{ssec:k01:setting}

Fix the couplings, a cube $\CJ$ and King's RG cube $\CK$ of \cref{def:s02:Kcube}. In this section and the next the index $k$ is the spacing index ($\eps_k=5^{-k}$, as throughout), and $j$ counts renormalisation steps, $0\le j\le k-J_*$ for a fixed $J_*$ (King's ``$J$'' of \cite[Thm.~3.1]{K1}); the block scale after $j$ steps is $\eta\defeq5^{j}\eps_k=5^{j-k}\le1$, and the unit lattice after $j$ steps is $\eta^{-1}(\CJ\cap\eta\Z^3)$. Since the side $2MJ$ of $\CJ$ is an integer, the faces of $\CJ$ lie on block planes at every scale.

\paragraph{Sources.} A real source is $s=(u,h)$ with $u=u_x\defeq\sum_ix_if_i$ a real combination of $2$-forms $f_i\in C^\infty_c(\R^3;\Lambda^2)$ sampled at the plaquette centres in the lattice's own frame, and $h\in C^\infty_c(\R^3)$ real; all supports lie at $\ell^\infty$-distance $>1$ from $\R^3\setminus\CJ$ (so $J\ge J_0$). Put $X(s)\defeq F(u)+\Wick{h}$, with the free Wick constant of \cref{def:s02:sources}; $X$ is real and linear in $s$. The unnormalised functional is
\[
\Zk(s)\ \defeq\ \int e^{-S}\,e^{X(s)}\ \prod dA\prod d^2\phi
\]
over the variables of $\CK$ at spacing $\eps_k$ (any fixed, $s$-independent normalisation of the measure; it cancels everywhere below), and $N_k(s)\defeq\Zk(s)/\Zk(0)=\avK{e^{X(s)}}$. At $h=0$ we write $a_k(x)\defeq N_k(u_x,0)$, which is the $a^K_{k,J}(x)$ of \cref{cor:s04:G2}.

\paragraph{King's flow on the cube.} The block-spin transformations of \cite[Sect.~2]{K1} are run on $\CK$ with block factor $5$: the photon field is averaged over the $5^{j}\eps_k$-blocks that tile $\XJ$, with the Neumann condition of \cref{def:s02:Kcube}; the scalar field is averaged with King's Dirichlet weight \cite[(3.8)--(3.9)]{K2} (in the corrected form in which the averaging blocks have their top layers on the Dirichlet planes) at the finest step and with the uniform weight above it. The Dirichlet planes of $\CK$ are $2MJ$ apart, a multiple of every block side, which is what makes this flow possible on $\CK$ and not on the comparison cube (\cref{rem:k01:cube}). Let $S^{(j)}_k(A_j,\phi_j;s)$ be King's effective action after $j$ steps \cite[(3.14)]{K1}, \cite[(3.12)]{K2}, with the source vertices \cite[(3.20)--(3.21)]{K1}; let $\chi_j$ be King's small-field indicator, \cite[(3.2)]{K1} times the face factor \cite[(5.41)]{K2}, a function of $(A_j,\phi_j)$ alone; and
\[
\Ik{j}(s)\ \defeq\ \int dA_j\,d\phi_j\ \chi_j\ e^{-S^{(j)}_k(A_j,\phi_j;s)} .
\]
The \emph{vacuum constant} $v_{k,j}\defeq S^{(j)}_k(0,0;0)$ is finite; it contains King's $E_0$, $E_1$, the normalisation constants $\ln Z_j$, $\ln Z_j(0)$ of the block-spin kernels and every field- and source-free graph. The \emph{normalised action} is $S^{(j)}_{N,k}(\cdot;s)\defeq S^{(j)}_k(\cdot;s)-v_{k,j}$ for every $s$ (the field-free, source-dependent part is kept), and $\Ihat{j}(s)\defeq\int\chi_je^{-S^{(j)}_{N,k}(\cdot;s)}=e^{v_{k,j}}\Ik{j}(s)$. The flow is a device of the proof: the state $\avK{\cdot}$ does not depend on it.

\paragraph{Diagonal pairs.} Fix $\alpha\in(0,\frac12)$ and put $j(k)\defeq\lfloor(1-\alpha)k\rfloor$, $\eta_k\defeq5^{j(k)-k}$; then $\eta_k\le5^{-\alpha k}$ and $k-j(k)\le\alpha k+1$. The \emph{diagonal pair} at $k$ consists of the $k$-model after $j(k)$ steps and the $(k+1)$-model after $j(k)+1$ steps. Both have the block scale $\eta_k$ ($5^{j+1}\eps_{k+1}=5^{j}\eps_k$), hence the same unit lattice and the same indicator $\chi_{j(k)}$; this is the situation of \cite[p.~656]{K1} (``the fields $A_k,\phi_k$ have the same bounds in each integral''). No relation between $j(k+1)$ and $j(k)+1$ is used anywhere.

\paragraph{Ba{\l}aban's decomposition.} In Ba{\l}aban's block-spin representation \cite{Ba1,Ba2}, $\Zk(s)=\int d\mu_{k,s}$ with $d\mu_{k,s}=e^{X(s)}d\mu_{k,0}$ a positive measure on the fields and on the block variables of the scales $\le j$, and $P=P_j\in[0,1]$ is the product of all small-field factors of those scales, a function of the block and fluctuation variables alone. The identity $\Zk(s)=\int P\,d\mu_{k,s}+\int(1-P)\,d\mu_{k,s}$ is exact; the first term is the \emph{small-field part}, the second the \emph{large-field remainder}.

\subsection{Two statements from the renormalisation-group analysis}\label{ssec:k01:open}

\begin{lemma}[\BALint: Ba{\l}aban's interaction-term bound at a nonzero background]\label{hyp:k01:BALint}
In the setting of \cite[Sect.~2]{Ba2} at block factor $L=5$, in unit-lattice form at the step $k$, $t=5^k\eps$: let $\Omega_A=B^k(\Lambda^{(k-1)\prime}_2)\supset\Omega_C=B^k(\Lambda^{(k)}_2)$ be the two propagator domains of \cite[Prop.~2.6]{Ba2}, let $\tilde B$ be the background field of \cite[p.~579]{Ba2}, regular on $\Omega_A$, let $G_k(\Omega,\tilde B)$ be the scalar propagator of the step on $\Omega$ in the background $\tilde B$, and write the interaction terms of the stage of \cite[p.~579]{Ba2}, localised in $B^k(\Lambda^{(k)}_7)$, as a sum of per-label full renormalised classes with scalar lines $G_k(\Omega_A,\tilde B)$, external scalar fields $E_A=a_kG_k(\Omega_A,\tilde B)f$ with $\norm f_\infty\le Cp\lambda(t)^{-1/4}$, and other external fields of norms $\le Cp^a\lambda(t)^{-b}$ (\cref{ssec:a07:dom}). Replace $G_k(\Omega_A,\tilde B)$ by $G_k(\Omega_C,B^{(k+1),\eta})$ in every scalar line and in every $E_A$. For every $\kappa>0$ there are $t_0(\kappa)>0$ and $C_\kappa<\infty$, depending on $\kappa$, on Ba{\l}aban's parameters $(R,\varrho,b_0,\tilde p,\bar n,\alpha_0)$ and on the couplings only, such that for $t\le t_0(\kappa)$ the sum over all localisations of the absolute values of all terms containing at least one difference of propagators or of external fields is at most $C_\kappa\,t^{\kappa}\,|\Lambda^{(k)}_7|$. The constants do not depend on $k$, on the domains or on $\tilde B$, and $\tilde B=0$ is included.
\end{lemma}
The proof is given in \cref{ssec:a07:dom} (\cref{thm:a07:dom}). This is the interaction-term half of Ba{\l}aban's domain-change estimate \cite[Prop.~2.6]{Ba2}, at every background. The terms with exactly one difference line, grouped line by line, carry their counterterm partners but not the complete renormalised classes of the cut graph, and a lone divergent two-photon-leg subgraph is not small per graph; the proof expands every line and every external field in the difference and regroups the terms exactly into contractions of \emph{full} classes of the cut components, to which the renormalisation bound \cite[Prop.~1, (1.33)]{Ba3} applies with constants free of the domains and of the background. Every use of the upper stability bound \cite[(1.3)]{Ba2} and of \cite[Thm.~3.1, (3.4)]{K1} at block factor $5$ passes through this statement.

\begin{lemma}[\RENorb: orbit regrouping of cut renormalised families]\label{hyp:k01:RENorb}
Let $\cR=\sum_{\Gamma\in\cG}c_\Gamma E(\Gamma)$ be a renormalised family of \cite[Sects.~1--3]{Ba3} with $\cG$ a finite union of renormalised classes (\cref{def:a07:class}), in which the objects (lines or external legs) of certain kinds are decomposed as $o=o'-\Delta o$ with symmetric kernels $\Delta o$, and expand multilinearly, $\cR=\sum_\tau2^{-|S_L|}c_\Gamma\prod_{o\in S}(-\Delta o)\prod_{o\notin S}o'$ over the oriented cut terms $\tau=(\Gamma,S,\theta)$ of \cref{def:a07:cut}. Call two cut terms equivalent when one arises from the other by finitely many of Ba{\l}aban's three completions of \cite[pp.~420, 429--432]{Ba3} (the counterterm partner; the two-vector-leg class; the permutation of the external scalar legs of a divergent $\phi\phi A$ subgraph) applied at divergent subgraphs containing no line of $S$, the objects of $S$ carried along to the image slots. Then: \textup{(i)} the equivalence classes (orbits) partition the cut terms, and every move keeps the label multiset and the number and kinds of the objects of $S$; \textup{(ii)} a connected line set of $\Gamma$ avoiding $S$ is a subgraph of the cut graph $\Gamma\setminus S$ with the same vertices, slots, slot types, order and degree \cite[(2.1)--(2.3)]{Ba3}, and cutting commutes with the moves, so that the cut graphs of one orbit form exactly one renormalised class of the cut graph, the product of the classes of its connected components; \textup{(iii)} at fixed labels, after the sum over the separable indices of the cut kernels, the sum of the cut terms of an orbit $O$ is $(-1)^{|S|}2^{-|S_L|}r_O^{-1}c_*\prod_KE(\Cl(K);\text{its fields})$, the product over the components $K$ of the class sums with the cut kernels as external fields, $c_*$ a constant and $r_O\ge1$ the number of relabellings of identical $S$-objects preserving the class; \textup{(iv)} the sum over all orbits is the full expansion; \textup{(v)} if $\cG$ is a per-label full class, the cut terms with a given cut datum (components with labels and orders, stubs and fields at each component, incidence of the $S$-objects) correspond $s_\delta$-to-one to the tuples of members of the per-label full classes of the components, members without internal line and of order $0$ included; each such group is a union of orbits and is closed under the exchange of the scalar legs of every two-scalar-leg or $\phi\phi A$ subgraph.
\end{lemma}
The proof is given in \cref{ssec:a07:orb} (\cref{lem:a07:orb,lem:a07:full}). This is Ba{\l}aban's regrouping step in the proof of \cite[Prop.~1, p.~432]{Ba3} (``if we remove a line in any graph of the class, then for the obtained graph we can find in this class all graphs necessary to form a new renormalized class''), made explicit for decomposed objects, with the sharper conclusion that one orbit gives exactly one class of the cut graph. It is needed wherever the renormalisation bound is applied to a family some of whose lines have been cut, in particular in the comparison of two propagator domains and in the telescoping of the cube's large-field remainders; at a nonzero background, where the bound of \cite[Prop.~1]{Ba3} is available for per-label full classes, it is used in the form \textup{(v)}.

\begin{remark}[the cube]\label{rem:k01:cube}
The propositions below are stated for King's RG cube $\CK$, on which King's Dirichlet weight is compatible with the image group of the two Dirichlet planes (they are $2MJ$ apart). On the comparison cube $\NDa$ the planes are $2MJ+\eps_k$ apart, no partition into consecutive blocks is invariant under the image group, and the flow does not exist; this is why \cref{hyp:s03:EC} is posed on King's convention, and why the fixed-cube limits of the comparison family are reached only through \cref{cor:s04:G2}. The cube of the existence theorem is $\CK$ itself (\cref{rem:s03:cube-identification}).
\end{remark}

\subsection{The inputs at zero scalar source}\label{ssec:k01:inputs}

\begin{proposition}[\Hrel: the small-field representation on King's cube, zero scalar source]\label{prop:k01:H0rel}
Fix $J\ge J_0$, a family $f_1,\dots,f_N$ and $x\in\R^N$, and let $s\in\{u_x,0\}$. There are $\eta_1>0$, $k_0$ and $C,c,\sigma>0$, depending on the couplings, on $J$, on $f$ and $x$, and on nothing else (not on $k$, $j$), such that for all $k\ge k_0$ and all $j$ with $J_*\le j\le k-J_*$ and $\eta=5^{j-k}\le\eta_1$:
\begin{equation}\label{eq:k01:H0rel}
e^{-C\eta^{\sigma}}\,\Ik{j}(s)\ \le\ \Zk(s)\ \le\ e^{C\eta^{\sigma}}\,\Ik{j}(s)+\rho(\eta)\,\Zk(0),\qquad \rho(\eta)\defeq e^{-c\,p(\eta)^2+C},
\end{equation}
where $p(\eta)=b_0(1+\ln\eta^{-1})^{p}$, $p\ge1$, is Ba{\l}aban's small-field threshold. Moreover, at $s=0$ the two terms of the upper line are the two parts of Ba{\l}aban's decomposition: $\int P\,d\mu_{k,0}\le e^{C\eta^\sigma}\Ik{j}(0)$ and $\int(1-P)\,d\mu_{k,0}\le\rho(\eta)\Zk(0)$ \textup{(\Rzero)}.
\end{proposition}
The proof is given in Appendix~\ref{app:a01_H0rel}. This is \cite[Thm.~3.1]{K1} on the cube, which \cite[Sect.~5.2]{K2} asserts (``the proof follows by modifying suitably the material in [Ba~1] and [Ba~2]''), in the form in which the remainder is measured against $\Zk(0)$ rather than against an absolute constant, and with constants uniform in the step $j$ (the form the diagonal pairs need, since $\eta_k\to0$). The lower line is the weaker of the two readings of the sign in \cite[(3.3)]{K1}. No lower bound on $\Zk(0)$ appears. Of the statements of \cref{ssec:k01:open}, \BALint\ (through Ba{\l}aban's torus bounds) and \RENorb\ (through the comparison of propagator domains and the telescoping of the large-field remainders) enter.

\begin{proposition}[\Kpos: one-step bound on the normalised action, zero scalar source]\label{prop:k01:Kpos}
With $J$, $f$, $x$ and $s\in\{u_x,0\}$ as above there are $\gamma\in(0,1)$, $\beta,\sigma>0$ and $C_3<\infty$, depending on the couplings, $J$, $f$, $x$ only, such that for all $k\ge k_0$, all $J_*\le j\le k-J_*$ with $\eta=5^{j-k}\le\eta_1$, and all $(A_j,\phi_j)$ in the support of $\chi_j$,
\begin{equation}\label{eq:k01:Kpos}
\big|S^{(j)}_{N,k}(A_j,\phi_j;s)-S^{(j+1)}_{N,k+1}(A_j,\phi_j;s)\big|\ \le\ C_3\big(5^{-\gamma j}\eta^{-\beta}+\eta^{\sigma}\big).
\end{equation}
\end{proposition}
The proof is given in Appendix~\ref{app:a02_Kpos}. This is \cite[Thm.~3.4]{K1} at $n=1$ on the cube, for the field- and source-dependent part of the effective action: the two members of a diagonal pair live on one unit lattice with one indicator $\chi_j$, and the statement compares them there. It is assembled from six pieces: the photon quadratic form (an even-restriction identity reduces the cube to the doubled torus, where \cite[Prop.~3.10]{K1} applies); the scalar quadratic form at zero background (the Dirichlet symbol) and with background (a two-leg tree graph); the background determinant (one-loop graphs with a remainder, \cite[(3.94)--(3.98)]{K1}, of which the middle line of (3.98) is false as printed and is replaced); the graphs with curvature vertices; the source-free field-dependent graphs away from the faces and near them (the photon face terms and the scalar boundary remainders of \cite[pp.~331--332]{K2}); and the block-vertex classes. For the curvature-vertex and the source-free field-dependent classes, the localisations whose tree length $R$ exceeds the small-field threshold $p$ are bounded by cutting the localised graph into clusters of mutual distance $>3\sqrt d$ and reducing each cluster in its own box of side $O(p)$, with the cut lines as external lines of both clusters (\cref{cor:a02:DEGall}(b) and \cref{lem:a02:SC}; the zone-B tails by \cref{thm:a02:Bsc}(a) and \cref{thm:a02:BVP}(d)). The vacuum constant $v_{k,j}$ is not part of the statement.

\begin{lemma}[bounds for free]\label{lem:k01:free}
Put $\beta_f\defeq C_0\sum_i\fournorm{f_i}^2$ with $C_0$ as in \cref{lem:s02:E}. For all $k\ge0$, $J\ge J_0$ and $x\in\R^N$,
\[
e^{-\beta_f|x|^2}\ \le\ a_k(x)\ \le\ e^{\beta_f|x|^2}.
\]
\end{lemma}
\begin{proof}
By \cref{thm:s05:CK}(a), $a_k(x)\le\exp(\norm{u_x}^2_{2,k})$, and $\norm{u_x}_{2,k}\le\sum_i|x_i|\norm{f_i}_{2,k}\le|x|(\sum_i\norm{f_i}^2_{2,k})^{1/2}\le|x|\,\beta_f^{1/2}$ by \cref{lem:s02:E}(E3). The lower bound is $a_k(x)a_k(-x)\ge1$, \cref{lem:s02:alg}(e).
\end{proof}

\subsection{The Cauchy property in the spacing}\label{ssec:k01:CA}

\begin{theorem}[the Cauchy property, ratio form]\label{thm:k01:CA}
Fix $J\ge J_0$, $f$, $x\in\R^N$, and assume \eqref{eq:k01:H0rel} and \eqref{eq:k01:Kpos} for $s\in\{u_x,0\}$ with the constants of \cref{prop:k01:H0rel,prop:k01:Kpos}. Choose $\alpha\in(0,\min\{\frac12,\gamma/(\gamma+\beta)\})$ and the diagonal pairs $j(k)=\lfloor(1-\alpha)k\rfloor$, $\eta_k=5^{j(k)-k}$. Put
\[
b_k\defeq C\eta_k^{\sigma},\qquad \theta_k\defeq C_3\big(5^{-\gamma j(k)}\eta_k^{-\beta}+\eta_k^{\sigma}\big),\qquad \rho_k\defeq\rho(\eta_k),\qquad r_k(0)\defeq2\rho_k,\qquad r_k(u_x)\defeq2\rho_ke^{\beta_f|x|^2}.
\]
Then there is $k_1$ such that for all $k\ge k_1$
\begin{equation}\label{eq:k01:CA}
\big|\ln a_{k+1}(x)-\ln a_k(x)\big|\ \le\ 4b_k+2\theta_k+r_k(u_x)+r_k(0).
\end{equation}
The right side is summable in $k$, and its tail from $k$ on is at most $C'5^{-\kappa k}$ for some $\kappa>0$. Hence $a_k(x)$ converges, and its limit lies in $[e^{-\beta_f|x|^2},e^{\beta_f|x|^2}]$.
\end{theorem}

\begin{proof}
Fix $k\ge k_0$ so large that $j(k)\ge J_*$, $\eta_k\le\eta_1$ and $r_k(u_x)\le1$ (possible by Step~6), and put $j\defeq j(k)$, $b\defeq b_k$, $\theta\defeq\theta_k$. For $s\in\{u_x,0\}$ write $Z\defeq\Zk(s)$, $I\defeq\Ik{j}(s)$, $\hat I\defeq\Ihat{j}(s)$ for the $k$-model after $j$ steps, and $Z'\defeq Z_{k+1}(s)$, $I'\defeq I^{(j+1)}_{k+1}(s)$, $\hat I'\defeq\hat I^{(j+1)}_{k+1}(s)$ for the $(k+1)$-model after $j+1$ steps. Both members have the block scale $\eta_k$ and the indicator $\chi_j$, and \eqref{eq:k01:H0rel} applies to both, since $j+1\le(k+1)-J_*$; their remainders are $R\defeq\rho_k\Zk(0)$ and $R'\defeq\rho_kZ_{k+1}(0)$.

\emph{Step 1 (the remainders relative to $Z$).} For $s=0$, $R/Z=\rho_k=r_k(0)/2$ and $R'/Z'=\rho_k$. For $s=u_x$, $Z=a_k(x)\Zk(0)\ge e^{-\beta_f|x|^2}\Zk(0)$ by \cref{lem:k01:free}, so $R/Z\le\rho_ke^{\beta_f|x|^2}=r_k(u_x)/2$, and likewise $R'/Z'\le r_k(u_x)/2$. In both cases $R/Z,\ R'/Z'\le r_k(s)/2\le\frac12$. No lower bound on $\Zk(0)$ and no bound on $\Zk(u_x)$ beyond \cref{lem:k01:free} is used.

\emph{Step 2 (one model).} From $e^{-b}I\le Z\le e^{b}I+R$: $I\le e^{b}Z$ and $I\ge e^{-b}(Z-R)=e^{-b}Z(1-R/Z)$. So $I\in[e^{-b}Z(1-R/Z),\,e^{b}Z]$, and in the same way $I'\in[e^{-b}Z'(1-R'/Z'),\,e^{b}Z']$. In particular $I,I'>0$.

\emph{Step 3 (across the pair).} On the support of $\chi_j$, \eqref{eq:k01:Kpos} gives $|S^{(j)}_{N,k}(\cdot;s)-S^{(j+1)}_{N,k+1}(\cdot;s)|\le\theta$. The two integrals $\hat I$, $\hat I'$ run over the same variables with the same $\chi_j$ and have non-negative integrands, so $\hat I'/\hat I\in[e^{-\theta},e^{\theta}]$. Since $I=e^{-v_{k,j}}\hat I$ and $I'=e^{-v_{k+1,j+1}}\hat I'$,
\[
I'/I\ \in\ e^{\Delta}\,[e^{-\theta},e^{\theta}],\qquad \Delta\defeq v_{k,j}-v_{k+1,j+1},
\]
and $\Delta$ does not depend on $s$: the vacuum constant is the same number in the numerator and in the denominator of $a_k$.

\emph{Step 4 (ratio of partition functions).} By Steps 2 and 3,
\[
\frac{Z'}{Z}=\frac{Z'}{I'}\cdot\frac{I'}{I}\cdot\frac{I}{Z}\ \in\ e^{\Delta}\Big[e^{-2b-\theta}\big(1-\tfrac RZ\big),\ e^{2b+\theta}\big(1-\tfrac{R'}{Z'}\big)^{-1}\Big].
\]
With $|\ln(1-t)|\le2t$ for $0\le t\le\frac12$ and Step~1, $|\ln(1-R/Z)|,\,|\ln(1-R'/Z')|\le r_k(s)$, hence
\[
\big|\ln(Z'/Z)-\Delta\big|\ \le\ 2b+\theta+r_k(s).
\]

\emph{Step 5 (ratio of ratios).} $\ln\dfrac{a_{k+1}(x)}{a_k(x)}=\Big[\ln\dfrac{Z_{k+1}(u_x)}{\Zk(u_x)}-\Delta\Big]-\Big[\ln\dfrac{Z_{k+1}(0)}{\Zk(0)}-\Delta\Big]$, and Step~4 at $s=u_x$ and at $s=0$ gives \eqref{eq:k01:CA}. (If the constants of \eqref{eq:k01:H0rel}, \eqref{eq:k01:Kpos} differ between the two sources, take the larger ones.)

\emph{Step 6 (summability).} Since $k-j(k)\ge\alpha k$, $\eta_k\le5^{-\alpha k}$, so $b_k$ and $C_3\eta_k^\sigma$ decay geometrically. Since $j(k)\ge(1-\alpha)k-1$ and $k-j(k)\le\alpha k+1$,
\[
5^{-\gamma j(k)}\eta_k^{-\beta}\ \le\ 5^{\gamma+\beta}\,5^{-(\gamma(1-\alpha)-\beta\alpha)k},
\]
which decays geometrically because $\alpha<\gamma/(\gamma+\beta)$. Finally $p(\eta_k)^2\ge b_0^2(1+\alpha k\ln5)^{2}\ge c_1k^2$ for $p\ge1$, so $\rho_k\le e^{-cc_1k^2+C}$; thus $r_k(u_x)\le1$ for large $k$, $\sum_kr_k(s)<\infty$, and the tail of $\sum_k\rho_k$ is eventually below any geometric tail. The series in \eqref{eq:k01:CA} converges with a geometric tail, so $(\ln a_k(x))_k$ is Cauchy, $a_k(x)$ converges, and the limit lies in the closed interval of \cref{lem:k01:free}.
\end{proof}

\begin{remark}[what the argument does not use]\label{rem:k01:notused}
(i) No convergence, and no bound, of the vacuum constants $v_{k,j}$ (King's $E_0$, $E_1$, the kernel normalisations, the face vacuum energies) enters: they cancel in Step~3. King's printed route \cite[p.~329]{K2} goes through the convergence of $\Zk$ itself, which requires exactly those constants to converge. (ii) No upper stability bound, no lower bound on $\Zk(0)$ and no bound on $\Zk(u_x)$ beyond \cref{lem:k01:free} is used; the remainder of \eqref{eq:k01:H0rel} is compared with $\Zk(0)$ only. (iii) No reflection positivity, no Ward identity, and no subsequence. (iv) The all-$n$ form of \cite[(3.13)]{K1} (comparing the $k$-model with the $(k+n)$-model at a fixed stopping scale $\eta$ for every $n$) also gives the Cauchy property, provided the constants of \eqref{eq:k01:H0rel} and \eqref{eq:k01:Kpos} are taken at fixed $\eta$ with no growth in $k$: then $|\ln a_{k+n}-\ln a_k|\le4b(\eta)+2\theta_k(\eta)+4\rho(\eta)(e^{\beta_f|x|^2}+1)$ for every $n$, and one lets $k\to\infty$ and then $\eta\to0$. The diagonal form above uses $n=1$ only, at the price of constants uniform in $j$; the two forms need different readings of the inputs, and neither implies the other.
\end{remark}

\begin{theorem}[King's existence theorem on the cube: the curvature field]\label{thm:k01:EC}
For every finite family $f_1,\dots,f_N$ of real $2$-forms in $C^\infty_c(\R^3;\Lambda^2)$, every $J\ge J_0(f)$ and every $x\in\R^N$ the limit
\[
\ZKJ(x)=\lim_{k\to\infty}\avK{e^{\sum_ix_iF(f_i)}}
\]
exists and lies in $[e^{-\beta_f|x|^2},e^{\beta_f|x|^2}]\subset(0,\infty)$, with $\beta_f$ as in \cref{lem:k01:free}; the convergence is geometric. That is, \cref{hyp:s03:EC} holds. Consequently $\ZKJ$ and $\ZKinf$ exist locally uniformly on $\C^N$ and are entire, with the bounds and the monotonicity of \cref{thm:s05:CK}(b).
\end{theorem}

\begin{proof}
\Cref{prop:k01:H0rel,prop:k01:Kpos} supply the hypotheses of \cref{thm:k01:CA} at every $x$; the theorem gives the limit and its bounds; \cref{thm:s05:CK}(b) (that is, \cref{cor:s04:G2}(ii)) gives the rest. The constants $k_0$, $\eta_1$ and the rates may depend on $x$ and on $f$; existence needs no uniformity in $x$.
\end{proof}

\begin{remark}[the comparison family]\label{rem:k01:compfamily}
\Cref{thm:k01:EC} says nothing about the fixed-cube limits of the comparison family $\NDa$ (\cref{rem:k01:cube}); for that family \cref{cor:s04:G2}(ii)($\gamma$) squeezes the limits inferior and superior at fixed $J$ between $\ZKJ$ and $Z^K_{J+1}$, and \cref{cor:s04:G2}(iii) identifies every subsequential limit, followed by $J\to\infty$, with $\ZKinf$. At fixed $J$ every inequality available for the comparison family bounds its oscillation in $k$ only by the difference of two distinct limits, which tends to zero only as $J\to\infty$; a Cauchy statement at fixed $J$ for that family would need the small-field analysis at its faces, where no flow of the kind used above exists.
\end{remark}
